\documentclass[11pt]{article}
\usepackage[margin=1in]{geometry}
\usepackage{amsmath,amssymb,amsthm,mathtools,hyperref,longtable}
\newtheorem{theorem}{Theorem}
\newtheorem{proposition}[theorem]{Proposition}
\newtheorem{lemma}[theorem]{Lemma}
\newtheorem{corollary}[theorem]{Corollary}
\theoremstyle{definition}\newtheorem{definition}[theorem]{Definition}
\newcommand{\dd}{d'd''}
\newcommand{\MA}{\operatorname{MA}}
\newcommand{\trop}{\operatorname{trop}}
\newcommand{\an}{\mathrm{an}}
\newcommand{\R}{\mathbb R}
\newcommand{\Z}{\mathbb Z}
\title{A period obstruction and bounded analytic Monge--Amp\`ere constructions\newline on an explicit compact analytic torus\newline Partial research toward P16.7}
\author{Unverified research draft}
\date{8 October 2026}
\begin{document}
\maketitle
\begin{abstract}
We construct a smooth proper analytic surface as an explicit lattice quotient,
prove an exact obstruction to positive integral period forms, and put a positive
real CLD form on the quotient. We construct genuine mixed CLD products on a
specified bounded continuous class, with mass and comparison. A bounded local
perturbation has a Monge--Amp\`ere atom of mass one outside the canonical
skeleton. We also solve every product-measure equation within a separable
skeletal class, with explicit energy, compactness, and one-dimensional
orthogonality. These results do not resolve P16.7. In particular the bridge
from the period obstruction to exclusion of every algebraization is not
certified under the requested whole-primary-proof audit, and no full analytic
finite-energy operator or all-support existence theorem is obtained.
\end{abstract}

\section{Verdict and retained scope}
The original target asks for an intrinsic finite-energy equation theory on
genuinely nonalgebraizable analytic spaces, including arbitrary analytic
support and all required compactness, envelope, first variation, continuity,
and comparison inputs. This draft proves the narrower results stated below.
It does not claim original completion, and its candidate flag is false.
No historical closure, reviewer decision, imported paper, or supplied program
is used as a correctness premise. No result on P16.1--P16.6 is counted as
progress here.

The field in all analytic constructions is
\[
 K=\mathbb C((s)),\qquad |s|=e^{-1/8},\qquad t=s^8,\qquad v=-\log|\cdot|.
\]
Thus $K$ is complete and discretely, nontrivially valued, with characteristic
and residue characteristic zero. Completeness follows directly by stabilization
of the coefficient of each Laurent power in a Cauchy sequence. Its residue
field is $\mathbb C$. We use good strictly analytic Berkovich spaces and
ordinary open sets. The constructed space is smooth, separated, proper,
boundaryless, geometrically connected, and pure of dimension two. The datum
is a real differential form, not an ample algebraic or formal line bundle.
All products here are CLD products. No identification with classical curve-test
psh, weak-smooth psh, residue-nef model psh, or a BFJ finite-energy class is
asserted.

\section{The quotient and the exact period calculation}
Put $T=(\mathbb G_m^2)^{\an}_K$, with coordinates $z_1,z_2$, and let
\[
 \lambda_1=(t,2),\qquad \lambda_2=(3,t),\qquad
 \lambda(n)=\lambda_1^{n_1}\lambda_2^{n_2}
       =(t^{n_1}3^{n_2},2^{n_1}t^{n_2}).
\]
The group $\Lambda=\{\lambda(n):n\in\Z^2\}$ acts by multiplication.
Set $x_i=-\log|z_i|$. Then
\[
 \trop(\lambda(n)z)=\trop(z)+n.                         \tag{2.1}
\]

\begin{proposition}\label{quotient}
The quotient $X=T/\Lambda$ exists as a good strictly analytic space and has
the geometric properties stated in Section 1. The quotient map is an ordinary
topological covering and a local analytic isomorphism. There are a continuous
map $\tau:X\to S=\R^2/\Z^2$ and a continuous closed section $i:S\to X$,
given by the coordinatewise Gauss norms.
\end{proposition}
\begin{proof}
For a real open rectangle $B$ with side lengths less than one, the open set
$U_B=\trop^{-1}(B)$ is disjoint from all its nonidentity $\Lambda$-translates.
Indeed a common point would give two elements of $B$ differing by a nonzero
integer vector. Equation (2.1) also shows that the action has no stabilizers.
The quotient of the underlying set is covered by the images of the $U_B$.
On an overlap the identifications are restrictions of multiplication by
$\lambda(n)$. They satisfy the cocycle rule because
$\lambda(n+m)=\lambda(n)\lambda(m)$. Gluing these open analytic spaces
therefore constructs $X$ and its analytic sheaf. The quotient topology makes
each of these maps an open embedding. This is an open analytic atlas, not a
claim based only on a G-cover.

For compact subsets $C,D\subset T$, only finitely many $n$ satisfy
$\lambda(n)C\cap D\ne\varnothing$: their integer vectors lie in the bounded
set $\trop(D)-\trop(C)$. As $T$ is locally compact and Hausdorff, two distinct
orbits can be separated by shrinking relatively compact neighborhoods and
then separating their finitely many potentially intersecting translates.
Hence $X$ is Hausdorff. Analytic separatedness follows as well: the inverse
image of its diagonal in $T\times_KT$ is the locally finite union of the
closed graphs of these translations; in the disjoint-sheet charts it is the
usual diagonal closed immersion. Thus the diagonal of $X$ is a closed
immersion.

The compact polyannulus $P=\trop^{-1}([0,1]^2)$ surjects onto $X$: translate
each tropical coordinate into $[0,1]$. It is strictly affinoid, since its
coordinate bounds are $|t|\le |z_j|\le1$. Thus $X$ is compact. It is locally
isomorphic to an open subset of the algebraic torus analytification, so is
smooth, pure of dimension two, good, strictly analytic, and boundaryless.
For goodness choose an atlas of rectangles with endpoints on the grid
$(1/8)\Z$ and side length $1/2$. Such rectangles cover $\R^2$ by their
interiors. Choose such a rectangle so each coordinate of the point lies
in its central half. Every point then has a strict closed polyannulus
neighborhood inside that sheet, with grid endpoints and interior containing that
point. Compactness, separatedness and absence of relative boundary give
properness in the Berkovich definition.

The polyannuli exhausting $T$ have integral Laurent-series algebras and are
connected: a clopen decomposition would give a nontrivial idempotent by
gluing the functions zero and one, impossible in an integral algebra. The
nested exhaustion is connected, hence so are $T$ and $X$. The same argument
works after any complete ground-field extension, since the polyannulus
Laurent-series algebras remain integral. This gives geometric connectedness.

For $x\in\R^2$ the rule
\[
 \eta_x\Big(\sum_m a_m z^m\Big)=\max_m |a_m|e^{-m\cdot x}
\]
defines a multiplicative Gauss norm. Multiplicativity follows, for example,
by passing to the associated graded Laurent polynomial ring, which is an
integral domain. It extends to convergent Laurent series. For polynomials
its dependence on $x$ is continuous, so $x\mapsto\eta_x$ is continuous in
the topology defined by evaluations. Multiplication by $\lambda(n)$ sends
$\eta_x$ to $\eta_{x+n}$. These norms descend to $i:S\to X$; the tropical
map descends to $\tau$ and $\tau i=\mathrm{id}$. Since $S$ is compact and
$X$ is Hausdorff, $i$ has closed image and is a homeomorphism onto it.
\end{proof}

\begin{definition}
A positive integral period form for the displayed lattice means a homomorphism
$c:\Lambda\to\Z^2$, interpreted as characters $z\mapsto z^{c(\lambda)}$,
such that
\[
 \lambda_j^{c(\lambda_i)}=\lambda_i^{c(\lambda_j)}
 \quad\text{and}\quad
 B_{ij}=v\big(\lambda_j^{c(\lambda_i)}\big)
 \text{ is positive definite}.
                                                               \tag{2.2}
\]
This is an explicitly defined period condition; the proposition below does
not identify it with analytic ampleness.
\end{definition}

\begin{proposition}\label{period}
Every homomorphism satisfying the symmetry in (2.2) has
$c(\lambda_1)=(0,k)$ and $c(\lambda_2)=(k,0)$ for some integer $k$.
Consequently no positive integral period form exists. The same obstruction
holds for rational character coefficients after clearing denominators, and
persists after every field extension containing $K$.
\end{proposition}
\begin{proof}
Write $c(\lambda_1)=(a,b)$ and $c(\lambda_2)=(c,d)$. The only nontrivial
symmetry equation is
\[
 3^a t^b=t^c2^d.
\]
Its valuation gives $b=c$. It then gives $3^a=2^d$. Unique factorization of
positive rational numbers implies $a=d=0$, including negative exponents
after clearing denominators. Therefore
\[
 B=\begin{pmatrix}0&k\\k&0\end{pmatrix}.
\]
If $k\ne0$ this has both a positive and a negative eigenvalue; if $k=0$ it
is zero. Neither sign convention for $B$ can make it positive definite.
Clearing a common denominator proves the rational assertion. All identities
used are identities between elements already in $K$, and remain exactly
the same in any extension field.
\end{proof}

\paragraph{Uncertified algebraization bridge.}
The published Riemann characterization of analytic tori and descent of an
analytic group law on a proper algebraization are the expected bridge from
Proposition \ref{period} to exclusion of every algebraization. It would be
incorrect to substitute the absence of a chosen equivariant polarization for
that bridge. The available complete CLD source, Appendix B.6, states a
classification of line bundles and an ampleness criterion but refers its
essential proofs to L\"utkebohmert, Section 6.2 and Theorem 6.4.4.
The full cited proofs were not retrieved. Fresnel--van der Put's Chapter 6
abstract authenticates the characterization but is not its proof. No theorem
of nonalgebraizability of the underlying $X$ is claimed in this draft under
the user's stronger source-audit rule. In particular the geometric example
does not yet meet the ``genuinely nonalgebraizable'' requirement as a certified
result. The gap concerns all algebraizations, including proper nonprojective
schemes and algebraic spaces, not just equivariant presentations.

\section{A positive real datum and genuine bounded mixed products}
Use the ordinary CLD calculus on $X$. Put
\[
 q(x)=\tfrac12(x_1^2+x_2^2),\qquad
 \theta=\dd q=\sum_{j=1}^2 d'x_j\wedge d''x_j.             \tag{3.1}
\]
This descends because $q(x+n)-q(x)=n\cdot x+|n|^2/2$ has zero CLD Hessian.
On each sheet the affine terms are real linear combinations of logarithms
of invertible analytic coordinates. The datum is therefore a globally
defined closed positive real $(1,1)$-form. Positivity here is positivity of
CLD forms; no assertion of strict positivity in every extra analytic direction
is made.

Our positive volume convention is
\[
 dV=d'x_1\wedge d''x_1\wedge d'x_2\wedge d''x_2,
 \qquad (\dd F)^2=2\det(D^2F)dV.                         \tag{3.2}
\]
The identity moment on the torus has analytic degree one, and the tangent
lattice is exactly $\Z^2$, hence saturated. Its canonical calibration is
one. By the definition of CLD integration in this chart,
\[
 \theta^2=2i_*(dx_1dx_2),\qquad \int_X\theta^2=2.          \tag{3.3}
\]
Lebesgue measure in (3.3) is on the unit square modulo translations and has
mass one. Thus the raw mass is two, not a probability normalization.

\begin{definition}
Let $\mathcal C_\theta$ be the real continuous functions $u$ on $X$ for which
$q+u$, on every sufficiently small sheet, is locally a uniform limit of
ordinary CLD-smooth functions with positive Hessian. This is an intrinsic
class relative to $\theta$: another lift changes $q$ by a smooth function
with zero Hessian. It is a bounded class because $X$ is compact. It is not
declared to be the full analytic psh class or its finite-energy domain.
\end{definition}

\begin{lemma}\label{boundedBT}
For $u,v\in\mathcal C_\theta$ there is a genuine positive symmetric mixed
CLD measure $\alpha_u\wedge\alpha_v$, where
$\alpha_u=\theta+\dd[u]$. It is local, has mass two, and is continuous
under uniform convergence within this class. In particular it is continuous
along decreasing sequences whose limit lies in $\mathcal C_\theta$.
Writing $\MA_\theta(u)=\alpha_u^2$, one has the comparison inequality
\[
 \MA_\theta(v)(\{u<v\})\le\MA_\theta(u)(\{u<v\}).        \tag{3.4}
\]
\end{lemma}
\begin{proof}
We describe the local construction, including the continuity input. If $T$
is a positive closed current and $f$ is a continuous locally uniformly
approximable psh function, define $\dd f\wedge T=\dd(fT)$.
Multiplication of a positive current by a continuous function is legitimate
because positive currents have order zero. On any fixed compact set their
pairings with a test form are bounded by pairings with a dominating positive
test form times the supremum norm of that function.
More explicitly, if $T_j\to T$ are positive and $f_j\to f$ locally uniformly,
choose a smooth $g$ with $\|g-f\|\le\varepsilon$ on a compact neighborhood
of the test support. The decomposition
\[
 f_jT_j-fT=(f_j-g)T_j+g(T_j-T)+(g-f)T
\]
and positive domination bound the limit superior of its pairing by a constant
times $\varepsilon$. Therefore $f_jT_j\to fT$. This is CLD Lemma 8.5.1,
whose complete proof is preserved and read. The initial smooth case is
$\dd(fT)=\dd f\wedge T$, positive and closed. Passing to locally uniform
limits proves positivity and closedness. This argument also supplies
continuity of the recursive products by induction. For smooth functions
they are ordinary exterior products. For two uniformly approximable
functions, approximate both smoothly; the symmetric smooth products converge
to both recursive orders, proving symmetry. These are precisely the CLD
Bedford--Taylor products, not an inverse measure construction.

On a sheet put $a=q$, $f=a+u$, $g=a+v$. The local product
$\dd(f\dd g)$ is positive. Adding to either potential a smooth function
with zero Hessian does not change it, by the smooth closed-current Leibniz
rule and symmetry. Hence the local products glue on $X$. Equivalently,
globally they are
\[
 \alpha_u\wedge\alpha_v
       =\theta\wedge\alpha_v+\dd(u\alpha_v).             \tag{3.5}
\]
This identity follows locally from $\dd(a\alpha_v)=\theta\wedge\alpha_v$.
Its last term has zero integral: its pairing with the smooth compactly
supported constant function one is its pairing with $\dd1=0$.
Likewise $\theta\wedge\alpha_v=\theta^2+\dd(v\theta)$ has the same
integral as $\theta^2$. This proves mass two and does not use an algebraic
line bundle. Local uniform continuity proves current convergence, hence weak
measure convergence by positive mass bounds and density of smooth tests.
For a decreasing sequence with continuous limit on compact $X$, Dini's
theorem gives uniform convergence. This does not prove continuity toward an
unbounded finite-energy limit.

The class is stable under maximum: on each sheet approximate $a+u$ and
$a+v$ smoothly and apply a smooth convex nondecreasing regularization of
maximum. The chain rule gives positive Hessian, and the regularizations
converge locally uniformly to $a+\max(u,v)$. Locality of (3.5) holds on
ordinary open sets by restriction of currents.
For $\varepsilon>0$ set $w=\max(u,v-\varepsilon)$. Locality and mass give
\[
 \MA_\theta(v)(\{u<v-\varepsilon\})
 \le2-\MA_\theta(u)(\{u>v-\varepsilon\})
 =\MA_\theta(u)(\{u\le v-\varepsilon\}).
\]
Let $\varepsilon\downarrow0$ and use continuity from below of finite
measures on both increasing sets. This proves (3.4).
\end{proof}

\section{A real analytic branch perturbation}
The following calculation uses the actual analytic quotient, not only its
abstract tropical lattice. Fix $r=|t|=e^{-1}$ and set
\[
 G(y)=\tfrac12y^2-\tfrac12y+\tfrac1{12}\quad (0\le y\le1),
 \qquad h_0=-G\quad\text{periodically},
 \qquad u_0=h_0(x_1)+h_0(x_2).                            \tag{4.1}
\]
Near $y=0$ the convex function $y^2/2+h_0(y)$ is $|y|/2-1/12$.
In particular $u_0\in\mathcal C_\theta$ and
$\MA_\theta(u_0)=2\delta_{i(0)}$, also as proved in Section 5.
Let $V\subset X$ be the sheet with $|x_i|<1/4$. Its coordinates are
unambiguously the original $z_i$. Define there
\[
 b_i=\log\max(1,|z_i|),\quad
 \ell_i=\log\max(r,|z_i-1|),\quad d_i=\ell_i-b_i,\quad
 D=\max(d_1,d_2).                                       \tag{4.2}
\]
Extend $D$ by zero outside $V$, and put $U=u_0+D$.

\begin{theorem}\label{branch}
The function $U$ is bounded continuous and belongs to $\mathcal C_\theta$.
It does not factor through $\tau$. Its genuine CLD measure has raw mass two
and has an atom of mass exactly one at the Gauss point $\xi_r$ of the closed
polydisc $|z_1-1|,|z_2-1|\le r$. This point is outside $i(S)$.
The statement does not specify the remaining mass-one part of the measure.
\end{theorem}
\begin{proof}
The ultrametric inequality gives $d_i\le0$. If $|z_i|\ne1$, then
$|z_i-1|=\max(1,|z_i|)$ and $d_i=0$. Thus $D$ can be nonzero only
when $x_1=x_2=0$. Its support is contained in the compact fiber
$\tau^{-1}(0)$, contained in the interior of $V$. The functions in (4.2)
are continuous even where $z_i-1$ vanishes because $r>0$. This proves
continuity of the extension and boundedness.

On $V$ a local potential for $\alpha_{u_0}$ is
\[
 F=q+u_0=b_1+b_2-\tfrac12\log|z_1z_2|-\tfrac16.
\]
The local potential for $\alpha_U$ is
\[
 F+D=\max\{\ell_1+b_2,b_1+\ell_2\}
           -\tfrac12\log|z_1z_2|-\tfrac16.              \tag{4.3}
\]
Each $b_i$ and $\ell_i$ is a local uniform limit of smooth psh functions.
For $\ell_i$, near a zero of $z_i-1$ it is constant on the open set
$|z_i-1|<r$; away from a zero it is the maximum of a constant and a
logarithm of an invertible analytic function. Smooth convex regularized
maximum proves the assertion in both cases and on their overlaps. The
logarithmic term in (4.3) has zero Hessian. Sums, and then regularized
maximum of the two sums, show that (4.3) is locally uniformly approximable
psh. Outside the support of $D$ the original local potential works. Hence
$U\in\mathcal C_\theta$. Lemma \ref{boundedBT} gives positivity and mass two.

Inside the ordinary open polydisc
$W=\{|z_1-1|<1,\ |z_2-1|<1\}\subset V$ one has $|z_i|=1$.
With $w_i=z_i-1$, (4.3) becomes
\[
 H=\log\max(r,|w_1|,|w_2|)-\tfrac16.                   \tag{4.4}
\]
Near $\xi_r$ both $w_i$ are invertible. The coordinates
$y_i=\log|w_i|$ identify this neighborhood with an open part of the
two-dimensional coordinate torus. Its tangent lattice is $\Z^2$ and its
analytic degree is one, since translation by $(1,1)$ is an analytic
isomorphism. The identity-moment skeleton there has exactly one point
over $(\log r,\log r)$, namely $\xi_r$.

For completeness compute its mass rather than rely on a truncated max
formula. Translate $y_i$ by $\log r$ and regularize
$\max(0,y_1,y_2)$ by
\[
 H_\varepsilon(y)=\varepsilon\log(1+e^{y_1/\varepsilon}
                                      +e^{y_2/\varepsilon}).
\]
Its gradient is the diffeomorphism from $\R^2$ onto
$\{(p_1,p_2):p_i>0,\ p_1+p_2<1\}$; the inverse is
$y_i=\varepsilon\log(p_i/(1-p_1-p_2))$. Therefore change of variables
gives $\int_{\R^2}\det D^2H_\varepsilon\,dy=1/2$.
Since $H_\varepsilon(y)=\varepsilon H_1(y/\varepsilon)$, these
determinant measures, of mass $1/2$, concentrate at zero: outside any
fixed ball their mass is the tail of the integrable determinant of $H_1$
outside a ball of radius tending to infinity. Formula (3.2) then gives
$2\det D^2H_\varepsilon\,dy\to\delta_0$.
The regularizations converge locally uniformly, so Lemma \ref{boundedBT}
and the identity torus calibration identify the analytic measure on a
neighborhood of $\xi_r$ with $\delta_{\xi_r}$. In fact the same calculation
on $W$ shows that the restriction of the measure to $W$ is this one atom.
Near $w_i=0$ any inactive logarithm can be omitted, and if fewer than two
independent logarithmic coordinates are active the square of the Hessian
is zero; this justifies covering also the coordinate-zero loci in $W$.

The point $\xi_r$ has $|z_i|=1$ but $|z_i-1|=r<1$.
The only point of $i(S)$ with $|z_i|=1$ for both $i$ is $i(0)$, where
the Gauss norm of $z_i-1$ is one. Hence $\xi_r\notin i(S)$.
Finally $D(i(0))=0$ and $D(\xi_r)=\log r=-1$, while both have tropical
coordinate zero. Thus $U$ cannot factor through $\tau$.
\end{proof}

\section{A completely explicit restricted equation and its energy}
We give all quantifiers of this restricted theorem. Let $C=\R/\Z$, with
Lebesgue probability $dy$, and let
\[
 \mathcal H=\{h\in C^0(C):\rho_h=dy+h''\ \text{is a positive distribution}\}.
                                                               \tag{5.1}
\]
Derivatives are distributional. The mass of $\rho_h$ is one by testing
$h''$ against the constant one. Equivalently $y^2/2+h(y)$ is convex
on every interval of its lifted domain.

\begin{lemma}\label{circle}
For each Radon probability $\nu$ on $C$, the unique mean-zero solution
of $h''=\nu-dy$ is
\[
 h_\nu(y)=-\int_C G(y-a)\,d\nu(a).                       \tag{5.2}
\]
It has $-1/12\le h_\nu\le1/24$, is $1/2$-Lipschitz, and has
oscillation at most $1/8$. Every member of $\mathcal H$ is of this form
up to an additive constant. The normalized class $\sup h=0$ is compact
in uniform topology.
\end{lemma}
\begin{proof}
The polynomial defining $G$ in (4.1) has equal endpoint values, mean zero,
second derivative one in $(0,1)$, and derivative jump $-1$ at zero.
Integration by parts against a smooth periodic test therefore gives
$G''=dy-\delta_0$. Convolution proves (5.2) distributionally.
The range of $G$ is $[-1/24,1/12]$ and its Lipschitz constant is $1/2$;
averaging its translates proves the asserted bounds and continuity.
If two solutions have the same second derivative their difference has
zero second derivative, hence is affine as a distribution on the line.
Periodicity makes it constant. This proves both uniqueness and the
description of $\mathcal H$.
After $\sup h=0$ the functions lie in $[-1/8,0]$ with common Lipschitz
constant $1/2$. Arzel\`a--Ascoli supplies subsequential uniform limits.
For a nonnegative smooth periodic test, the positive distribution
$dy+h''$ remains positive under uniform limits. Thus the normalized
class is closed, proving compactness.
\end{proof}

Define the separable analytic class
\[
 \mathcal S_\theta=\{u=h_1(\tau_1)+h_2(\tau_2):h_i\in\mathcal H\}.
\]
Only the projections of $\tau$ are meant; $X$ need not be a product of
elliptic curves. Convex convolution of $y^2/2+h(y)$ on the line gives
locally uniform smooth convex approximations, so
$\mathcal S_\theta\subset\mathcal C_\theta$.

\begin{theorem}\label{separable}
For $u=h_1+h_2$ and $v=k_1+k_2$ in this class the genuine mixed analytic
measure is
\[
 \alpha_u\wedge\alpha_v
  =i_*\big(\rho_{h_1}\otimes\rho_{k_2}
                   +\rho_{k_1}\otimes\rho_{h_2}\big).     \tag{5.3}
\]
Thus every measure $2i_*(\nu_1\otimes\nu_2)$ for arbitrary Radon
probabilities $\nu_j$ on $C$ has a solution in $\mathcal S_\theta$,
unique in this class modulo constants, namely (5.2) in each coordinate.
No assertion of uniqueness among all analytic psh functions is made.
\end{theorem}
\begin{proof}
For smooth $h_i,k_i$ the Hessian forms are diagonal, with entries
$1+h_i''$ and $1+k_i''$. Their wedge is the sum of the two cross
products in (5.3), by ordinary exterior algebra. The identity torus
moment gives support $i(S)$ and multiplicity one, so this is an equality
of analytic measures, not only of tropical pushforwards.
Convolve each $h_i$ periodically with a smooth nonnegative approximate
identity. Its curvature measure is $\rho_{h_i}$ convolved with that
identity, hence positive and of mass one. The convolutions converge
uniformly. Lemma \ref{boundedBT} proves convergence of the analytic
products. The product probabilities converge weakly as well: this first
holds for products of continuous tests, then for all continuous tests
by uniform approximation of functions on $C^2$ by finite sums of
product functions. Hence (5.3) holds for all members of the class.
For $u=v$ it becomes $2i_*(\rho_{h_1}\otimes\rho_{h_2})$.
Taking probability marginals recovers each $\rho_{h_i}$. Lemma
\ref{circle} therefore gives the solution and uniqueness modulo constants.
\end{proof}

\subsection{Energy, dual energy, and first variation within this class}
Every $h\in\mathcal H$ is Lipschitz, so the following is finite:
\[
 e(h)=\int_C h\,dy-\tfrac12\int_C|h'|^2\,dy,
 \qquad E_{\rm sep}(u)=2(e(h_1)+e(h_2)).                 \tag{5.4}
\]
This is independent of the decomposition: changing it adds a constant to
one $h_i$ and subtracts it from the other. Adding $c$ to $u$ adds $2c$
to $E_{\rm sep}$, consistent with raw mass two.
It agrees on this class with the mixed CLD energy
\[
 E(u)=\tfrac13\int_Xu(\theta^2+\theta\wedge\alpha_u+\alpha_u^2).
                                                               \tag{5.5}
\]
Indeed (5.3) makes the three summands equal respectively to
$2\sum_i\int h_i\,dy$, $\sum_i(\int h_i\,dy+\int h_i\,d\rho_{h_i})$,
and $2\sum_i\int h_i\,d\rho_{h_i}$. Their average is
$\sum_i(\int h_i\,dy+\int h_i\,d\rho_{h_i})$, and periodic integration
by parts gives
$\int h_i\,d\rho_{h_i}=\int h_i\,dy-\int|h_i'|^2dy$.
For Lipschitz members this identity follows by periodic convolution:
derivatives converge in $L^2$, and the uniform/weak convergence of the
function/curvature measures permits passage to the limit.

For $h,k\in\mathcal H$ the exact difference identity is
\[
 e(k)-e(h)=\int_C(k-h)\,d\rho_h
                  -\tfrac12\int_C|k'-h'|^2dy.           \tag{5.6}
\]
The same convolution argument proves the integration by parts used here.
Consequently $e$ is concave and strictly concave modulo constants.
Along every segment in $\mathcal S_\theta$, differentiation of (5.4)
and (5.6) gives
\[
 \frac d{da}E_{\rm sep}(u+a(v-u))
          =\int_X(v-u)\,\MA_\theta(u+a(v-u)).            \tag{5.7}
\]
The identity is proved only for these directions, not for arbitrary
analytic envelope variations.

For a positive analytic Radon measure $\mu$ of mass two define its
\emph{restricted} dual energy by
\[
 E^*_{\rm sep}(\mu)=
 \sup_{u\in\mathcal S_\theta}\big(E_{\rm sep}(u)-\int_Xu\,d\mu\big).
                                                               \tag{5.8}
\]
Constants cancel in (5.8). Normalize both $\sup h_i=0$.
Then $-1/4\le u\le0$ and $E_{\rm sep}(u)\le0$, so (5.8) lies
between zero and $1/2$ for every such measure $\mu$. In particular this
condition is far too weak to be substituted for the full analytic
finite-energy measure condition.
For $\mu=2i_*(\nu_1\otimes\nu_2)$, (5.6) implies that (5.2) maximizes
$E_{\rm sep}-\int u\,d\mu$, uniquely modulo constants, and its
measure is the prescribed $\mu$ by Theorem \ref{separable}.

More generally let $\nu_i=\tfrac12(\tau_i)_*\mu$. The same exact
quadratic identity shows that the restricted maximizer is still (5.2),
but its analytic measure is
$2i_*(\nu_1\otimes\nu_2)$, which need not equal $\mu$.
For example take $y$ to be the rigid point represented by $z_1=z_2=1$
and $\mu=2\delta_y$. Then $E^*_{\rm sep}(\mu)<\infty$ and its
restricted maximizer is $u_0$, with measure $2\delta_{i(0)}\ne2\delta_y$.
In this example $E_{\rm sep}(u_0)=-1/6$, $u_0(y)=-1/6$, and the exact
restricted dual energy is $E^*_{\rm sep}(2\delta_y)=1/6$ by (5.6).
This is an actual analytic counterexample to the proposed use of the
restricted class and its dual energy for all-support existence. It is
not a counterexample to P16.7 with a correctly constructed full psh cone;
we do not claim this rigid-point measure has finite energy in that cone.

\subsection{One-dimensional envelopes and orthogonality, with proof}
For any continuous obstacle $a:C\to\R$ let
$P(a)=\sup\{h\in\mathcal H:h\le a\}$.
A constant below $\min a$ makes the set nonempty. The common Lipschitz
bound in Lemma \ref{circle} and the obstacle bound show that its pointwise
supremum is finite and $1/2$-Lipschitz. On the lifted line each
$y^2/2+h$ is convex, so their supremum is convex. Therefore $P(a)$ belongs
to $\mathcal H$, is below $a$, and is the maximum member of the set.
It follows directly that
\[
 \|P(a)-P(b)\|_\infty\le\|a-b\|_\infty.                \tag{5.9}
\]

On a small closed interval contained in $\{P(a)<a\}$ with positive
clearance from the obstacle, write $F(y)=y^2/2+P(a)(y)$. If $F$ is not
affine on a sufficiently short such interval, replace it there by its
endpoint chord. The chord lies above $F$, and its slope lies between
the outside one-sided slopes of $F$, so the replacement is convex.
Uniform continuity makes its increase smaller than the obstacle
clearance by choosing the interval short enough. Make the same
replacement in every integer translate; the transformation rule
$F(y+1)=F(y)+y+1/2$ is preserved. This contradicts maximality.
If $F''$ were nonzero somewhere in the open noncontact set, there would
be an arbitrarily short interval of this kind on which it is not affine.
Hence $F''=0$ there, and
\[
 \rho_{P(a)}\text{ is supported on }\{P(a)=a\},\qquad
 \int_C(a-P(a))\,d\rho_{P(a)}=0.                        \tag{5.10}
\]

For completeness, put $p=P(a)$ and $p_b=P(a+bv)$ for continuous $v$.
The difference identity (5.6), its reverse, and (5.10) give
\[
 b\int_C v\,d\rho_{p_b}
 \le e(p_b)-e(p)\le b\int_Cv\,d\rho_p.                 \tag{5.11}
\]
For the upper bound use $p=a$ $\rho_p$-almost everywhere and
$p_b\le a+bv$; for the lower bound use $p_b=a+bv$
$\rho_{p_b}$-almost everywhere and $p\le a$.
By (5.9), $p_b\to p$ uniformly as $b\to0$, and thus the
distributional curvature probabilities converge weakly. Dividing (5.11)
by $b$, reversing inequalities if $b<0$, proves
\[
 \left.\frac d{db}\right|_{b=0}e(P(a+bv))
                  =\int_C v\,d\rho_{P(a)}.              \tag{5.12}
\]
These are circle envelopes. They are not envelopes over all analytic
$\theta$-psh functions on $X$, even for a separable obstacle.

\section{Sources, audit boundaries, and unresolved quantifiers}
The full unchanged original task and all three manifests were read first.
The historical mathematical packets were inventoried as data; their results
are not used. The entire supplied CLD-v2 PDF and literal text, and the
supplied GR-v1 PDF and literal text, are preserved locally. Direct applied
CLD inputs and complete proofs read are: positive currents and their
order-zero structure in Section 7.2; continuous psh and convex regularized
maxima in Sections 8.1--8.2; local estimates in 8.4; Lemmas 8.5.1,
8.5.3, 8.5.5; the recursive construction and smooth-Hessian invariance
in Sections 8.6--8.7; Propositions 8.7.2 and 8.7.4; and the
multilinear dpsh construction in 8.8.5. The torus Gauss norms and
calibration definitions in Sections 12.1 and 14.3 were read with their
surrounding hypotheses and proofs. For our degree-one coordinate moments
the calibration follows directly from its definition. CLD PDF pages 209
and 213 were viewed to check the recursive-product and max conventions.
The mass-one max calculation in Section 4 is proved independently above.
The analytic-space applicability statement is CLD Theorem 15.7.1:
a topologically separated pure $n$-dimensional analytic space, with its
analytic boundary and canonical calibration, is a tropical space of level
$n$. Its full immediate proof and hypotheses were read; here $n=2$ and
the boundary is empty. The smooth-test separation, density, partitions,
and positive-test decomposition inputs in Sections 15.2, 6.2, 6.5 and
Appendix A.2 were also read with their immediate proofs. The proof above
uses only closed currents for recursion; it does not use the claim about
arbitrary positive continuous multiplication in Corollary 8.6.7.

Recursive published foundations are not silently granted: CLD uses its
earlier construction of smooth forms, partitions and smooth-test density,
analytic tropical calibrations, and standard Berkovich local geometry.
These are the foundational differential/analytic setting of this draft;
not every recursive published proof behind that setting was re-proved.
None of them supplies global analytic finite-energy compactness or
envelope regularity. A separate exact source-application ledger records
what was applied, what was only located, and what is missing.

The public search examined post-2020 analytic semipositivity and toric
pluripotential work, the 2026 balanced-polyhedral manuscript, and the
1 October 2026 Kaib--Werner abeloid-uniformization preprint. The last is
actual analytic geometric literature, not an analytic finite-energy
Monge--Amp\`ere theorem. GR's analytic PL results likewise are not such an
equation theorem. BMPS's polyhedral equations and its algebraic analytic
comparisons cannot be imported as an all-support theorem on this quotient.
Search coverage is finite and does not establish a universal absence of
literature. The complete proofs of these works are not claimed audited
where they are not applied to a theorem above. Some literature bodies
were available through the public web tool, while the inherited shell
proxy was unreachable and the two torus books failed full-body retrieval.
No proxy, authentication or CA configuration was changed.

The missing original requirements are: a certified proof excluding every
algebraization of the example; an intrinsic full analytic psh cone and
its full $\mathcal E^1$ domain; positive mixed products on that domain
with unbounded decreasing continuity and full mass; its correct dual
energy condition for arbitrary-support analytic measures; global
compactness and analytic envelopes; orthogonality and first variation
in all analytic directions; existence for every measure satisfying that
condition; and uniqueness modulo constants in the full domain.
The bounded comparison lemma is a real partial input, not a replacement
for all these missing quantifiers. No analytic counterexample satisfying
the whole original target's ambient requirements is asserted.

\section{Reproduction and computational evidence}
This file is standalone LaTeX. The complete companion proof is
\texttt{whole-proof.md}; the exact application and coverage records are
\texttt{source-application-ledger.md} and \texttt{scope-and-verdict.json}.
The authorized source-copy script is \texttt{scripts/prepare\_evidence.py}.
The bounded diagnostic is \texttt{scripts/diagnostics.py}; run it with
\texttt{python scripts/diagnostics.py}. It checks rational identities,
the exhaustive list of symmetric character matrices in a stated finite
box, and the factorial/simplex and Green-function normalizations.
The proofs above are universal arguments independent of those tests.
No supplied program was executed or imported, no agents or extra model
calls were created, and no external changes or messages were made.
There is no resource-heavy computation and no such computation was
launched. TeX compilation was not performed because the user reserves
compilation to ROOT. A lexical artifact check is recorded separately;
it is not a TeX compilation or mathematical certification.

\begin{thebibliography}{9}
\bibitem{CLD} A. Chambert-Loir and A. Ducros,
\emph{Formes diff\'erentielles r\'eelles et courants sur les espaces de
Berkovich}, supplied version 2, 25 July 2025.
\bibitem{GR} W. Gubler and J. Rabinoff,
\emph{On semipositive piecewise linear functions in non-archimedean
analytic geometry}, arXiv:2503.05513v1, 7 March 2025.
\bibitem{BMPS} A. M. Botero, E. Mazzon, L. Pille-Schneider,
\emph{Monge--Amp\`ere measures on balanced polyhedral spaces},
arXiv:2603.08664v1, March 2026.
\bibitem{KW} S. Kaib and A. Werner,
\emph{Abeloid varieties over higher rank valued fields: Raynaud
uniformization and N\'eron models}, arXiv:2610.01354v1, 1 October 2026.
\bibitem{FvdP} J. Fresnel and M. van der Put,
\emph{Rigid Analytic Geometry and Its Applications}, Chapter 6,
2004. Full immediate torus proofs not retrieved; not used as a proved
nonalgebraization premise.
\bibitem{L} W. L\"utkebohmert,
\emph{Rigid Geometry of Curves and Their Jacobians}, 2016.
Full cited immediate proofs not retrieved; not used as a proved
nonalgebraization premise.
\end{thebibliography}
\end{document}
